\documentclass[12pt]{article}

\usepackage[utf8]{inputenc}
\usepackage[spanish]{babel}
\usepackage[backend=biber]{biblatex}
\usepackage[letterpaper, margin=2.54cm]{geometry}
\usepackage{csquotes, multicol}
\setlength{\columnsep}{1cm}
\setlength{\parindent}{0cm}
\pagestyle{plain}
\addbibresource{referencias.bib}

\title{Investigación de Operaciones \\
	\vspace{0.4cm}
	\Large Una breve introducción
}
\date{04 de julio de 2026}
\author{Rosendo Camal}

\begin{document}

% Portada
\maketitle

\begin{multicols}{2}

% Introducción 
La \textit{Investigación de Operaciones} es una herramienta que ayuda a las organizaciones, como las empresas, a la toma de decisiones. Para comprender su finalidad es necesario comprender qué es, qué objetivos cumple y la manera en qué contribuye a ciertas organizaciones.
\newline

% Desarrollo

La \textbf{Investigación de Operaciones} es un arte y una ciencia que ayuda a tomar mejores decisiones a instituciones, empresas, gobiernos o a cualquier actor involucrado en una situación cuya manera de abordarlo sea priorizando la aplicación de la solución más óptima y considerando sus restricciones o circunstancias. Es un \textbf{arte} porque la naturaleza de las circunstancias y de las problemáticas enfrentadas conscientemente ocasiona que el proceso del modelado ---del problema--- no sea rígido, predecible, fácilmente estandarizable y replicable. Por lo que las metodologías heurísticas, metaheurísticas y el sentido empírico juegan un papel para reconocer, definir, comprender y resolver la crítica situación, aún cuando se empleen técnicas de análisis matemáticas. Por otro lado, llamarla \textbf{ciencia} es precisa dada la necesidad y presencia de cálculos, modelos matemáticos y algoritmos junto a la implementación del método científico. La elección correcta, o la más adecuada, de las técnicas matemáticas viene apoyada de forma sistémica por el método mencionado, a un proceso de iteración basado en el ensayo y error, y la ayuda del pesamiento matemático orientado a la solución de problemas.
\newline

Aunque definir en tan pocas palabras el campo de investigación de operaciones puede resultar una empresa arriesgada debido al peligro de dar a conocer superficialmente su complejidad y su amplitud. A pesar de ello, sin demasiado esfuerzo podemos apreciar que  el núcleo radica en la \textbf{resolución de problemas}. He aquí, por tanto, su objetivo fundamental: hallar la solución de ciertos problemas. Es trivial la similitud con el espíritu de las matemáticas ---puras y aplicadas---, pero aquí nos enfrentamos a la necesidad de acotar y colocar una línea divisoria en cuanto a lo qué es un problema para la \textit{investigación de operaciones}, qué se considera solucionar y cómo se solucionan aquellos planteamientos. Dada la magnitud de estas cuestiones, no son necesarias abordarlas en esta breve introducción, pero es adecuado conocer estos cuestionamientos.
\newline

Además del objetivo dado a conocer previamente, se pueden destacar otras que giran en torno al principal:

\begin{itemize}
	\item Resolver el problema de la mejor manera posible.
	\item Comprender el alcance del problema y sus consecuencias.
	\item Uso de heurísticas o técnicas matemáticas para lograr el fin deseado.
	\item Resolución, amortiguación o prevención de problemas en las operaciones de una organización.
	\item Mejor aprovechamiento de los recursos, es decir, una mejor administración.
\end{itemize}

Reconozco que lo anterior puede considerarse disperso y hasta críptico, sin embargo, esto es así ya que la \textit{investigación de operaciones} es el resultado de la intersección de las matemáticas aplicadas y la administración. Un dato, más curioso que relevante, de esta área del saber fue que surgió durante la \textbf{Segunda Guerra Mundial}, un episodio oscuro de la \textquote{avanzada civilización} del Siglo XX, donde numerosos científicos y militares tomaron la consigna de elaborar planes y estrategias aprovechando óptimamente los recursos bélicos disponibles (dinero, tiempo, soldados, armas, vehículos, alimentos, entre otros). Como consecuencia el nombre derivó del acto de \textquote{Investigar Operaciones [de origen militar]}.
\newline

La mayor contribución de la \textit{investigación de operaciones} a las organizaciones, en especial, a las empresas es permitir tener recursos o herramientas de origen, principalmente, matemático para dar fin a problemáticas reconocidas a través de los indicadores de la empresa o mediante la natural necesidad, deseo o ambición de optimizar recursos (maximizar ganancias, reducir costes).
\newline

	Para ilustrar la contribución de la investigación de operaciones a una empresa podemos plantear un caso sencillo e hipotético como el caso de unas \textquote{quejas en la línea de espera de un elevador}~\cite{taha}. En esta situación, el equipo de expertos estudiaron la coyuntura enfocándose en el comportamiento de los quejosos. Esto causó que eliminaran una ansiada implementación de un modelo sofisticado y complejo basado en la teoría de colas y optaran por colocar un espejo. Esta sencilla y económica idea se le ocurrió al psicológo del grupo, y en contra de lo que aparenta, eliminó al 100\% las molestias de los usuarios del elevador ya que preferían pasar el tiempo de espera mirándose entre ellos y viéndose en el espejo, que impacientarse y quejarse. En otras situaciones será necesario la implementación de modelos matemáticos, pero también podría ocurrir que no fuesen necesarios y al final, lo primordial es dar fin a un problema mediante el mejor uso de los recursos disponibles.
\newline

% Conclusión

En resumen, la \textit{Investigación de Operaciones} es una ciencia administrativa surgida durante la Segunda Guerra Mundial y que emplea técnicas de matemáticas aplicadas. Igualmente, se puede considerar un arte dado la naturaleza heurística, metaheurística y empírica de ella. Su finalidad mayor es hallar la solución a problemas específicos, y siendo una herramienta aprovechada por organizaciones como empresas, quienes se ven beneficiadas por el mejor uso de los recursos que disponen.

% Bibliografía

\nocite{*}
\printbibliography[heading=bibintoc]

\end{multicols}
\end{document}
